\chapter{State of the Art in Spaceborne Wildfire
Monitoring: History and Current
Gaps}\label{chapter-4-state-of-the-art-in-spaceborne-wildfire-monitoring--history-and-current-gaps}

\begin{quote}
\textbf{Writing Guide \& Methodological Focus:}

\begin{itemize}
\item
  \textbf{Objective:} Show how wildfire EO evolved, what systems exist
  today, and where they still miss or mission requirements.
\item
  \textbf{Method:} Tell the \textbf{history by epochs} → compile a
  \textbf{systems comparison table} → add a brief \textbf{modality
  overview} → close with a \textbf{gap matrix} vs requirements.
\item
  \textbf{Caution:} Tell a story of progress and remaining challenges,
  don\textquotesingle t just list missions. Describe capabilities and
  limits; do \textbf{not} design the constellation here.
\item
  \textbf{Checkpoint:} Does this chapter successfully demonstrate that,
  despite decades of progress, a specific and important capability gap
  still exists? Can a reader point to concrete shortfalls this work
  design must close (tempospatial resolution(GSD, latency), observation
  in core bands, all-weather reliability)?
\end{itemize}
\end{quote}

\section{Historical evolution}\label{41-historical-evolution}

Four short phases---Foundations (1970-1999), Global Monitoring
(1999--2012), Multi-sensor convergence (2013--2019), NewSpace \& rapid
analytics (2020--present)---ending with ``lessons learned.''

\subsection{The Foundational Era (1970s --
1990s)}\label{411-the-foundational-era-1970s----1990s}

\begin{itemize}
\item
  Cover the first thermal detections and post-fire mapping with Landsat.
\end{itemize}

\subsection{The Era of Global Monitoring (2000 --
2015)}\label{412-the-era-of-global-monitoring-2000----2015}

\begin{itemize}
\item
  Discuss MODIS, the impact of the open-data policy, and the rise of
  commercial VHR.
\item
  Introduce the Copernicus Sentinel missions.
\end{itemize}

\subsection{The NewSpace Revolution (2015 --
Present)}\label{413-the-newspace-revolution-2015----present}

\begin{itemize}
\item
  Cover VIIRS, commercial constellations (Planet, ICEYE), and
  purpose-built fire constellations (FireBIRD, OroraTech, FireSat).
\end{itemize}

\section{Comparative Synthesis EO Wildfire
systems}\label{42-comparative-synthesis-eo-wildfire-systems}

One table with columns: \textbf{Name \textbar{} First launch \textbar{}
Operational \textbar{} Application \textbar{} Region of interest
\textbar{} Orbit \textbar{} \# sats \textbar{} Spatial res \textbar{}
Temporal res \textbar{} VNIR \textbar{} SWIR \textbar{} MWIR \textbar{}
LWIR \textbar{} SAR \textbar{} FOV (deg) \textbar{} Swath @ nadir (km)
\textbar{} Max off-nadir (deg) \textbar{} Tip-and-cue (Y/N) \textbar{}
Slew/retask (typ.) \textbar{} Tasking autonomy (Y/N)};

conclude with the single-sentence takeaway that there is \textbf{no
regional, all-weather constellation} meeting \textless30 min sub-hour
cadence + \textless10m meter-class thermal over Portugal's daily fire
cycle.

\section{Modality overview (capabilities \&
limits)}\label{43-modality-overview-capabilities--limits}

In one paragraph each, summarize \textbf{GEO}, \textbf{SSO},
\textbf{RGTO}, \textbf{polar medium-res thermal}, \textbf{high-res
optical/thermal}, \textbf{SAR/microwave}---what each does well for fires
and where each misses for the regional Portuguese case.§2.6 (cadence,
GSD, weather/angle robustness, latency). In the methodology, a detailed
trade-off of these types of orbit will be done. This subchapter shall be
used to filter from the table other requirements for our modality

\section{Regional fit: Portugal vs seasonal
windows.}\label{44-regional-fit-portugal-vs-seasonal-windows}

How current systems actually cover Portugal's \textbf{seasonal} and
\textbf{core local-time} bands; note where QC (cloud/smoke/angle)
removes scenes.

\section{Latency, access, and tasking
realities.}\label{46-latency-access-and-tasking-realities}

Typical acquisition→alert latency ranges, downlink/ground constraints,
and practical tasking/retasking windows during peak fire hours. Connect
to the Portuguese architecture of wildfire management.

\section{Coverage geometry, swath, agility \& tip-and-cue
capability}\label{47-coverage-geometry-swath-agility--tip-and-cue-capability}

Show that wide \textbf{FOV} without \textbf{pointing agility} reduces
\textbf{effective swath} and increases complexity in optics, difficulty
the coverage in core fire daily cycle windows, and note the emerging
constellation with \textbf{two-layer layers tip-and-cue} (wide FOV
``tipper'' + agile narrow FOV ``cuer'') as a way to tile Portugal's
seasonal bands more efficiently decreasing hardware constraints.
\textbf{Output:} A requirement seed: \emph{support tip-and-cue within
core local-time bands, respecting maximum allowed view-zenith angle (θ)
and ≤{[}TBD{]} retask latency}.

\section{Synthesis, Mission Requirements Vs EO Wildfire
Capability-Gap
Matrix}\label{48--synthesis-mission-requirements-vs-eo-wildfire-capability-gap-matrix}

One-page matrix aligning \textbf{mission requirements} (AOI/ROIs,
seasonal \& daily fire windows, max revisit, GSD, θ-cap, latency,
transferability) against each EO Wildfire mission/modality, highlighting
\textbf{where gaps remain} and which are only partially covered. Check
if new requirements are needed.

